% Part: first-order-logic
% Chapter: natural-deduction
% Section: soundness

\documentclass[../../../include/open-logic-section]{subfiles}

\begin{document}

\iftag{FOL}
      {\olfileid{fol}{ntd}{sou}}
      {\olfileid{pl}{ntd}{sou}}

\olsection{યથાર્થતા}

\begin{explain}
સ્વાભાવિક નિગમન જેવી કોઈ !!^a{derivation} પદ્ધતિ વાસ્તવમાં ફલિત ન થતી
બાબતોને !!{derive} ન કરી શકે, તો તે \emph{યથાર્થ} છે. તેથી યથાર્થતા
એ !!{derivation} પદ્ધતિઓ માટે બાંયધરીવાળો સુરક્ષા-ગુણધર્મ છે. કયો
સાબિતી-સૈદ્ધાંતિક ગુણધર્મ વિચારીએ છીએ તેના આધારે, દાખલા તરીકે, આપણે
જાણવા માગીએ કે
\begin{enumerate}
\item દરેક !!{derivable} !!{sentence} \iftag{FOL}{પ્રામાણ્ય}{પુનરુક્તિ} છે;
\item કોઈ !!a{sentence} બીજાં અમુક વાક્યોમાંથી !!{derivable} હોય, તો તે
  તેમનું ફલિત પણ છે;
\item !!{sentence}sનો કોઈ ગણ અસુસંગત હોય, તો તે અસંતોષ્ય છે.
\end{enumerate}
!!a{derivation} પદ્ધતિ માટે આ મહત્વના ગુણધર્મો છે. તેમાંનો કોઈ સધાય
નહીં તો !!{derivation} પદ્ધતિ ખામીવાળી ગણાય---તે જરૂર કરતાં વધારે
!!{derive} કરશે. તેથી કોઈ !!{derivation} પદ્ધતિની યથાર્થતા સ્થાપવી
અત્યંત મહત્વની છે.
\end{explain}

\begin{thm}[યથાર્થતા]
\ollabel{thm:soundness}
જો $!A$ એ !!{undischarged} ધારણાઓ~$\Gamma$માંથી !!{derivable} હોય,
તો $\Gamma \Entails !A$.
\end{thm}

\begin{proof}
$!A$ની કોઈ !!a{derivation}ને $\delta$ કહો. $\delta$માંનાં અનુમાનોની
સંખ્યા પર આગમનથી આગળ વધીએ.

આગમનના આધાર માટે અનુમાનોની સંખ્યા~$0$ હોય ત્યારે દાવો બતાવીએ. આ
કિસ્સામાં $\delta$માં માત્ર એક !!{sentence}~$!A$, એટલે કે એક ધારણા,
હોય છે. એ ધારણા !!{undischarged} છે, કારણ કે ધારણાઓ માત્ર અનુમાનોથી
!!{discharged} થઈ શકે છે અને અહીં કોઈ અનુમાન નથી. તેથી સાબિતીની બધી
!!{undischarged} ધારણાઓને સંતોષતી કોઈ પણ
\iftag{FOL}{!!{structure}~$\Struct{M}$}{!!{valuation}~$\pAssign{v}$}
એ~$!A$ને પણ સંતોષે છે.

હવે આગમન-પગલું લો. માનો કે $\delta$માં~$n$ અનુમાનો છે. સૌથી નીચેના
અનુમાનનાં આધારવાક્યો એવી ઉપ-!!{derivation}s વડે !!{derive}d થયાં છે
જેમાં દરેકમાં~$n$ કરતાં ઓછાં અનુમાનો છે. આગમન-પરિકલ્પના પ્રમાણે સૌથી
નીચેના અનુમાનનાં આધારવાક્યો એ આધારવાક્યો પર સમાપ્ત થતી
ઉપ-!!{derivation}sની !!{undischarged} ધારણાઓમાંથી ફલિત થાય છે. આખી
સાબિતીની !!{undischarged} ધારણાઓમાંથી ફલિતવાક્ય~$!A$ ફલિત થાય છે તેમ
આપણે બતાવવું છે.

સૌથી નીચેના અનુમાનના પ્રકાર પ્રમાણે કિસ્સા પાડીએ. પહેલાં માત્ર એક
આધારવાક્ય ધરાવતાં સંભવિત અનુમાનો જોઈએ.

\begin{enumerate}
\item માનો કે છેલ્લું અનુમાન \Intro{\lnot} છે. !!{derivation}નું
  સ્વરૂપ આ છે:
  \begin{prooftree}
    \AxiomC{$\Gamma, \Discharge{!A}{n}$}
    \RightLabel{$\delta_1$}
    \DeduceC{$\lfalse$}
    \DischargeRule{\Intro{\lnot}}{n}
    \UnaryInfC{$\lnot !A$}
  \end{prooftree}
  આગમન-પરિકલ્પના મુજબ $\lfalse$ એ~$\delta_1$ની !!{undischarged}
  ધારણાઓ $\Gamma \cup \{!A\}$માંથી ફલિત થાય છે. કોઈ
  \iftag{FOL}{!!a{structure}~$\Struct{M}$}{!!a{valuation}~$\pAssign{v}$}
  લો. જો $\iftag{FOL}{\Sat{M}}{\pSat{v}}{\Gamma}$ હોય, તો
  $\iftag{FOL}{\Sat{M}}{\pSat{v}}{\lnot !A}$ છે તેમ બતાવવું છે.
  પ્રતિવાદ માટે માનો કે $\iftag{FOL}{\Sat{M}}{\pSat{v}}{\Gamma}$,
  પરંતુ $\iftag{FOL}{\Sat/{M}}{\pSat/{v}}{\lnot !A}$, એટલે કે
  $\iftag{FOL}{\Sat{M}}{\pSat{v}}{!A}$. ત્યારે
  $\iftag{FOL}{\Sat{M}}{\pSat{v}}{\Gamma \cup \{!A\}}$ થશે, જે આપણી
  આગમન-પરિકલ્પનાની વિરુદ્ધ છે. તેથી
  $\iftag{FOL}{\Sat{M}}{\pSat{v}}{\lnot !A}$.

\item છેલ્લું અનુમાન \Elim{\land} છે. તેનાં બે રૂપ છે: આધારવાક્ય
  $!A \land !B$માંથી $!A$ અથવા $!B$નું અનુમાન થઈ શકે છે. પહેલો
  કિસ્સો લો. !!{derivation}~$\delta$ આવી દેખાય છે:
  \begin{prooftree}
    \AxiomC{$\Gamma$}
    \RightLabel{$\delta_1$}
    \DeduceC{$!A \land !B$}
    \RightLabel{\Elim{\land}}
    \UnaryInfC{$!A$}
  \end{prooftree}
  આગમન-પરિકલ્પના મુજબ $!A \land !B$ એ~$\delta_1$ની
  !!{undischarged} ધારણાઓ~$\Gamma$માંથી ફલિત થાય છે. કોઈ
  !!a{structure}~\iftag{FOL}{$\Struct{M}$}{$\pAssign{v}$} લો. જો
  $\iftag{FOL}{\Sat{M}}{\pSat{v}}{\Gamma}$ હોય, તો
  $\iftag{FOL}{\Sat{M}}{\pSat{v}}{!A}$ છે તેમ બતાવવું છે. માનો કે
  $\iftag{FOL}{\Sat{M}}{\pSat{v}}{\Gamma}$. આગમન-પરિકલ્પના
  ($\Gamma \Entails !A \land !B$) મુજબ
  $\iftag{FOL}{\Sat{M}}{\pSat{v}}{!A \land !B}$. વ્યાખ્યા મુજબ
  $\iftag{FOL}{\Sat{M}}{\pSat{v}}{!A \land !B}$ હોય જો અને માત્ર જો
  $\iftag{FOL}{\Sat{M}}{\pSat{v}}{!A}$ અને
  $\iftag{FOL}{\Sat{M}}{\pSat{v}}{!B}$. ($!A \land !B$માંથી $!B$નું
  અનુમાન થતું હોય તે કિસ્સો એ જ રીતે સંભાળાય છે.)

\item છેલ્લું અનુમાન \Intro{\lor} છે. તેનાં બે રૂપ છે: આધારવાક્ય
  ~ $!A$ અથવા આધારવાક્ય~$!B$માંથી $!A \lor !B$નું અનુમાન થઈ શકે છે.
  પહેલો કિસ્સો લો. !!{derivation}નું સ્વરૂપ આ છે:
  \begin{prooftree}
    \AxiomC{$\Gamma$}
    \RightLabel{$\delta_1$}
    \DeduceC{$!A$}
    \RightLabel{\Intro{\lor}}
    \UnaryInfC{$!A \lor !B$}
  \end{prooftree}
  આગમન-પરિકલ્પના મુજબ $!A$ એ~$\delta_1$ની !!{undischarged}
  ધારણાઓ~$\Gamma$માંથી ફલિત થાય છે. કોઈ
  \iftag{FOL}{!!a{structure}~$\Struct{M}$}{!!a{valuation}~$\pAssign{v}$}
  લો. જો $\iftag{FOL}{\Sat{M}}{\pSat{v}}{\Gamma}$ હોય, તો
  $\iftag{FOL}{\Sat{M}}{\pSat{v}}{!A \lor !B}$ છે તેમ બતાવવું છે.
  માનો કે $\iftag{FOL}{\Sat{M}}{\pSat{v}}{\Gamma}$; ત્યારે આગમન-પરિકલ્પના
  $\Gamma \Entails !A$ મુજબ $\iftag{FOL}{\Sat{M}}{\pSat{v}}{!A}$.
  તેથી $\iftag{FOL}{\Sat{M}}{\pSat{v}}{!A \lor !B}$ પણ હોવું જ જોઈએ.
  ($!B$માંથી $!A \lor !B$નું અનુમાન થતું હોય તે કિસ્સો એ જ રીતે
  સંભાળાય છે.)

\item છેલ્લું અનુમાન \Intro{\lif} છે. ધારણા~$!A$ અને
  ફલિતવાક્ય~$!B$ ધરાવતી ઉપસાબિતીમાંથી $!A \lif !B$નું અનુમાન થાય છે,
  એટલે કે
  \begin{prooftree}
    \AxiomC{$\Gamma, \Discharge{!A}{n}$}
    \RightLabel{$\delta_1$}
    \DeduceC{$!B$}
    \DischargeRule{\Intro{\lif}}{n}
    \UnaryInfC{$!A \lif !B$}
  \end{prooftree}
  આગમન-પરિકલ્પના મુજબ $!B$ એ~$\delta_1$ની !!{undischarged}
  ધારણાઓમાંથી ફલિત થાય છે, એટલે કે $\Gamma \cup \{!A\} \Entails !B$.
  કોઈ \iftag{FOL}{!!a{structure}~$\Struct{M}$}{!!a{valuation}~$\pAssign{v}$}
  લો. $!A$ છેલ્લાં અનુમાન વખતે નિવૃત્ત થાય છે, તેથી~$\delta$ની
  !!{undischarged} ધારણાઓ માત્ર $\Gamma$ છે. એટલે આપણે
  $\Gamma \Entails !A \lif !B$ બતાવવું છે. પ્રતિવાદ માટે માનો કે કોઈ
  \iftag{FOL}{!!{structure}~$\Struct{M}$}{!!{valuation}~$\pAssign{v}$}
  માટે $\iftag{FOL}{\Sat{M}}{\pSat{v}}{\Gamma}$ પરંતુ
  $\iftag{FOL}{\Sat/{M}}{\pSat/{v}}{!A \lif !B}$. તેથી
  $\iftag{FOL}{\Sat{M}}{\pSat{v}}{!A}$ અને
  $\iftag{FOL}{\Sat/{M}}{\pSat/{v}}{!B}$. પરંતુ પરિકલ્પના મુજબ $!B$ એ
  $\Gamma \cup \{!A\}$નું ફલિત છે, એટલે
  $\iftag{FOL}{\Sat{M}}{\pSat{v}}{!B}$, જે વ્યાઘાત છે. તેથી
  $\Gamma \Entails !A \lif !B$.

\item છેલ્લું અનુમાન \FalseInt છે. અહીં $\delta$નો અંત આ છે:
  \begin{prooftree}
    \AxiomC{$\Gamma$}
    \RightLabel{$\delta_1$}
    \DeduceC{$\lfalse$}
    \RightLabel{\FalseInt}
    \UnaryInfC{$!A$}
  \end{prooftree}
  આગમન-પરિકલ્પના મુજબ $\Gamma \Entails \lfalse$. આપણે
  $\Gamma \Entails !A$ બતાવવું છે. નહિ એમ માનો; તો કોઈ
  ~ $\iftag{FOL}{\Struct{M}}{\pAssign{v}}$ માટે
  $\iftag{FOL}{\Sat{M}}{\pSat{v}}{\Gamma}$ અને
  $\iftag{FOL}{\Sat/{M}}{\pSat/{v}}{!A}$. પરંતુ હંમેશાં
  $\iftag{FOL}{\Sat/{M}}{\pSat/{v}}{\lfalse}$ હોય છે, તેથી તેનો અર્થ
  $\Gamma \Entails/ \lfalse$ થશે, જે આગમન-પરિકલ્પનાની વિરુદ્ધ છે.

\item છેલ્લું અનુમાન \FalseCl છે: અભ્યાસપ્રશ્ન.

\iftag{FOL}{
\item છેલ્લું અનુમાન \Intro{\lforall} છે. ત્યારે $\delta$નું સ્વરૂપ આ છે:
  \begin{prooftree}
    \AxiomC{$\Gamma$}
    \RightLabel{$\delta_1$}
    \DeduceC{$!A(a)$}
    \RightLabel{\Intro{\lforall}}
    \UnaryInfC{$\lforall[x][!A(x)]$}
  \end{prooftree}
  આગમન-પરિકલ્પના મુજબ આધારવાક્ય $!A(a)$ એ !!{undischarged}
  ધારણાઓ~$\Gamma$નું ફલિત છે. $\Sat{M}{\Gamma}$ હોય એવી કોઈ
  સંરચના~$\Struct{M}$ લો. આપણે $\Sat{M}{\lforall[x][!A(x)]}$ બતાવવું છે.
  $\lforall[x][!A(x)]$ !!a{sentence} હોવાથી તેનો અર્થ દરેક
  ચલ-નિયુક્તિ~$s$ માટે $\Sat{M}{!A(x)}[s]$ બતાવવો એવો થાય છે
  (\olref[syn][ass]{prop:sat-quant}). $\Gamma$માં માત્ર વાક્યો હોવાથી
  \olref[syn][sat]{defn:satisfaction} મુજબ દરેક $!B \in \Gamma$ માટે
  $\Sat{M}{!B}[s]$. $\Struct{M'}$ એ~$\Struct{M}$ જેવી જ હોય, ફક્ત
  $\Assign{a}{M'} = s(x)$ હોય. $a$ એ~$\Gamma$માં આવતું નથી, તેથી
  \olref[syn][ext]{cor:extensionality-sent} મુજબ $\Sat{M'}{\Gamma}$.
  $\Gamma \Entails !A(a)$ હોવાથી $\Sat{M'}{!A(a)}$. $!A(a)$
  !!a{sentence} હોવાથી \olref[syn][ass]{prop:sentence-sat-true} મુજબ
  $\Sat{M'}{!A(a)}[s]$. \olref[syn][ext]{prop:ext-formulas} મુજબ
  $\Sat{M'}{!A(x)}[s]$ હોય જો અને માત્ર જો $\Sat{M'}{!A(a)}$ (યાદ કરો
  કે $!A(a)$ એ માત્ર $\Subst{!A(x)}{a}{x}$ છે). તેથી
  $\Sat{M'}{!A(x)}[s]$. $a$ એ~$!A(x)$માં આવતું નથી, તેથી
  \olref[syn][ext]{prop:extensionality} મુજબ $\Sat{M}{!A(x)}[s]$.
  પરંતુ $s$ મનસ્વી ચલ-નિયુક્તિ હતી, તેથી
  $\Sat{M}{\lforall[x][!A(x)]}$.

\item છેલ્લું અનુમાન \Intro{\lexists} છે: અભ્યાસપ્રશ્ન.

\item છેલ્લું અનુમાન \Elim{\lforall} છે: અભ્યાસપ્રશ્ન.
}{}
\sourcecorrection{OLND-006}{સ્થિર અંગ્રેજી મૂળ સર્વપરિમાણકના
નિવારણ-કિસ્સાને અહીં એકમાત્ર વાર $\Elim{\forall}$ લખે છે; પ્રકરણના
નિયમ અને બાકીના બધા ઉપયોગ મુજબ તેને $\Elim{\lforall}$ કર્યું છે.}
\end{enumerate}

હવે અનેક આધારવાક્યો ધરાવતાં સંભવિત અનુમાનો જોઈએ:
\Elim{\lor}, \Intro{\land},
\iftag{FOL}{\Elim{\lif} અને \Elim{\lexists}}{\Elim{\lif}}.
\begin{enumerate}

\item છેલ્લું અનુમાન \Intro{\land} છે. આધારવાક્યો $!A$ અને $!B$માંથી
  $!A \land !B$નું અનુમાન થાય છે અને $\delta$નું સ્વરૂપ આ છે:
  \begin{prooftree}
    \AxiomC{$\Gamma_1$}
    \RightLabel{$\delta_1$}
    \DeduceC{$!A$}
    \AxiomC{$\Gamma_2$}
    \RightLabel{$\delta_2$}
    \DeduceC{$!B$}
    \RightLabel{\Intro{\land}}
    \BinaryInfC{$!A \land !B$}
  \end{prooftree}
  આગમન-પરિકલ્પના મુજબ $!A$ એ~$\delta_1$ની !!{undischarged}
  ધારણાઓ~$\Gamma_1$માંથી અને $!B$ એ~$\delta_2$ની !!{undischarged}
  ધારણાઓ~$\Gamma_2$માંથી ફલિત થાય છે. ~ $\delta$ની !!{undischarged}
  ધારણાઓ $\Gamma_1 \cup \Gamma_2$ છે, તેથી
  $\Gamma_1 \cup \Gamma_2 \Entails !A \land !B$ બતાવવું છે. કોઈ
  \iftag{FOL}{!!a{structure}~$\Struct{M}$}{!!a{valuation}~$\pAssign{v}$}
  લો જેને $\iftag{FOL}{\Sat{M}}{\pSat{v}}{\Gamma_1 \cup \Gamma_2}$.
  $\iftag{FOL}{\Sat{M}}{\pSat{v}}{\Gamma_1}$ હોવાથી
  $\Gamma_1 \Entails !A$ પ્રમાણે $\iftag{FOL}{\Sat{M}}{\pSat{v}}{!A}$;
  અને $\iftag{FOL}{\Sat{M}}{\pSat{v}}{\Gamma_2}$ હોવાથી
  $\Gamma_2 \Entails !B$ પ્રમાણે $\iftag{FOL}{\Sat{M}}{\pSat{v}}{!B}$.
  બંને સાથે $\iftag{FOL}{\Sat{M}}{\pSat{v}}{!A \land !B}$.

\item છેલ્લું અનુમાન \Elim{\lor} છે: અભ્યાસપ્રશ્ન.

\item છેલ્લું અનુમાન \Elim{\lif} છે. આધારવાક્યો $!A \lif !B$ અને
  ~ $!A$માંથી $!B$નું અનુમાન થાય છે. !!{derivation}~$\delta$ આવી દેખાય છે:
  \begin{prooftree}
    \AxiomC{$\Gamma_1$}
    \RightLabel{$\delta_1$}
    \DeduceC{$!A \lif !B$}
    \AxiomC{$\Gamma_2$}
    \RightLabel{$\delta_2$}
    \DeduceC{$!A$}
    \RightLabel{\Elim{\lif}}
    \BinaryInfC{$!B$}
  \end{prooftree}
  આગમન-પરિકલ્પના મુજબ $!A \lif !B$ એ~$\delta_1$ની
  !!{undischarged} ધારણાઓ~$\Gamma_1$માંથી અને $!A$ એ~$\delta_2$ની
  !!{undischarged} ધારણાઓ~$\Gamma_2$માંથી ફલિત થાય છે. કોઈ
  \iftag{FOL}{!!a{structure}~$\Struct{M}$}{!!a{valuation}~$\pAssign{v}$}
  લો. જો $\iftag{FOL}{\Sat{M}}{\pSat{v}}{\Gamma_1 \cup \Gamma_2}$ હોય,
  તો $\iftag{FOL}{\Sat{M}}{\pSat{v}}{!B}$ બતાવવું છે. માનો કે
  $\iftag{FOL}{\Sat{M}}{\pSat{v}}{\Gamma_1 \cup \Gamma_2}$. કારણ કે
  $\Gamma_1 \Entails !A \lif !B$, તેથી
  $\iftag{FOL}{\Sat{M}}{\pSat{v}}{!A \lif !B}$. કારણ કે
  $\Gamma_2 \Entails !A$, તેથી $\iftag{FOL}{\Sat{M}}{\pSat{v}}{!A}$.
  તેનો અર્થ $\iftag{FOL}{\Sat{M}}{\pSat{v}}{!B}$ થાય છે (કારણ કે જો
  $\iftag{FOL}{\Sat/{M}}{\pSat/{v}}{!B}$ હોય, તો
  $\iftag{FOL}{\Sat{M}}{\pSat{v}}{!A}$ હોવાથી
  $\iftag{FOL}{\Sat/{M}}{\pSat/{v}}{!A \lif !B}$ થાય, જે
  ~ $\iftag{FOL}{\Sat{M}}{\pSat{v}}{!A \lif !B}$ને વિરુદ્ધ છે).

\item છેલ્લું અનુમાન \Elim{\lnot} છે: અભ્યાસપ્રશ્ન.

\tagitem{FOL}{છેલ્લું અનુમાન \Elim{\lexists} છે: અભ્યાસપ્રશ્ન.}{}
\end{enumerate}
\end{proof}

\tagprob{FOL}
\begin{prob}
\olref[fol][ntd][sou]{thm:soundness}ની સાબિતી પૂરી કરો.
\end{prob}
\tagendprob

\tagprob{notFOL}
\begin{prob}
\olref[pl][ntd][sou]{thm:soundness}ની સાબિતી પૂરી કરો.
\end{prob}
\tagendprob

\begin{cor}
\ollabel{cor:weak-soundness}
જો $\Proves !A$ હોય, તો $!A$ \iftag{FOL}{પ્રામાણ્ય}{પુનરુક્તિ} છે.
\end{cor}

\begin{cor}
\ollabel{cor:consistency-soundness}
જો $\Gamma$ સંતોષ્ય હોય, તો તે સુસંગત છે.
\end{cor}

\begin{proof}
આપણે પ્રતિવિધાન સાબિત કરીએ. માનો કે $\Gamma$ સુસંગત નથી. ત્યારે
$\Gamma \Proves \lfalse$, એટલે કે~$\Gamma$માંની !!{undischarged}
ધારણાઓમાંથી $\lfalse$ની કોઈ !!a{derivation} છે. \olref{thm:soundness}
મુજબ $\Gamma$ને સંતોષતી કોઈ પણ
\iftag{FOL}{!!{structure}~$\Struct{M}$}{!!{valuation}~$\pAssign{v}$}
એ $\lfalse$ને સંતોષવી જોઈએ. દરેક
\iftag{FOL}{!!{structure}~$\Struct{M}$}{!!{valuation}~$\pAssign{v}$}
માટે $\iftag{FOL}{\Sat/{M}}{\pSat/{v}}{\lfalse}$ હોવાથી કોઈ
\iftag{FOL}{$\Struct{M}$}{$\pAssign{v}$} એ~$\Gamma$ને સંતોષી શકતી નથી;
એટલે કે $\Gamma$ સંતોષ્ય નથી.
\end{proof}

\end{document}
